% Derived from output/experiment2/formal_20260917_v1/derived/analysis.json.
% The source data and formal plan remain unchanged. Requires only standard LaTeX.
\begin{table*}[t]
\centering
\small
\setlength{\tabcolsep}{4pt}
\renewcommand{\arraystretch}{1.15}
\begin{tabular}{@{}p{0.13\textwidth}p{0.12\textwidth}p{0.36\textwidth}p{0.32\textwidth}@{}}
\hline
Problem & Rank check & Unscreened random features & Constructed basis \\
\hline
Cubic PINN & $Q,P$: 10/10 & $\max E/B=0.972$; $\min(B_{256}-E_{256})=1.34\times10^{-3}$ & not applicable \\
$p=4$ Ritz & $D$: 10/10 & $\max\Delta_{\rm fun}=6.64\times10^{-4}$; $\max D_{\rm end}=0.235$; $\max\rho_4=1.20$ & $\Delta_{\rm fun}=1.63\times10^{-8}$; $D_{\rm end}=1.23\times10^{-8}$; $\max\rho_4=3.00\times10^{-7}$ \\
Allen--Cahn & $V$: 10/10 & $\max\rho_{\rm AC}=0.146$; $8/10$ final residuals exceed their initial values & $\rho_{\rm AC}=4.48\times10^{-8}$ (numerical stop) \\
\hline
\end{tabular}
\caption{Experiment 2, independent-quadrature diagnostics. Rank counts use 10 fixed feature seeds. All unscreened random-feature and constructed $p=4$ trajectories reached the 50,000-step budget. Constructed Allen--Cahn stopped at step 9,579 due to numerical step underflow, not the gradient threshold. Stationarity residuals use an independent 256-point rule. The constructed basis is illustrative, not representative of random features.}
\label{tab:experiment2}
\end{table*}
